\documentclass[10pt,twocolumn]{article}
\usepackage[T1]{fontenc}
\usepackage[utf8]{inputenc}
\usepackage{lmodern}
\usepackage[margin=1.8cm,top=2.4cm,bottom=2cm,columnsep=0.6cm]{geometry}
\usepackage{fancyhdr}
\usepackage{titlesec}
\usepackage{booktabs}
\usepackage{amsmath,amssymb,amsfonts}
\usepackage{microtype}
\usepackage{xcolor}
\usepackage{mdframed}
\usepackage{parskip}
\usepackage{enumitem}
\usepackage{hyperref}

\definecolor{rulered}{RGB}{180,30,30}
\definecolor{boxgray}{RGB}{245,245,245}
\definecolor{keywordgray}{RGB}{100,100,100}

\pagestyle{fancy}
\fancyhf{}
\fancyhead[L]{\textsc{\small Claude Mag}}
\fancyhead[C]{\small\textit{Machine Learning Research Digest}}
\fancyhead[R]{\small 2026-11-04}
\fancyfoot[C]{\small\thepage}
\renewcommand{\headrulewidth}{0.8pt}

\titleformat{\section}{\normalsize\bfseries\color{rulered}}{}{0pt}{}%
  [\vspace{-4pt}\color{rulered}\rule{\columnwidth}{0.4pt}]
\titlespacing{\section}{0pt}{8pt}{4pt}

\hypersetup{colorlinks=true,urlcolor=rulered,linkcolor=black}

\begin{document}
\setlength{\parindent}{0pt}
\setlength{\parskip}{4pt}

{\small\color{keywordgray}\textbf{Keywords:}
  zeroth-order optimization \textbullet{}
  reinforcement learning for LLMs \textbullet{}
  noise amplification \textbullet{}
  memory-efficient fine-tuning}\par\vspace{4pt}

{\LARGE\bfseries\raggedright
  The Cost of Mismatch}\par\vspace{4pt}

{\large\itshape\raggedright
  Noise Amplification in Zeroth-Order Reinforcement Learning
  Scales as $1/\mu^2$; Matched-Data BOLT Cancels It}\par\vspace{6pt}

{\small\textbf{By} Lianmin Chen, Junbin Qiu, Chenxing Wei, Yao Shu,
  Kun He}\par\vspace{2pt}

{\small\color{keywordgray}arXiv:2611.03260 \textbullet{}
  cs.LG \textbullet{} NeurIPS 2026 Poster}
\par\vspace{2pt}\noindent{\color{rulered}\rule{\columnwidth}{0.6pt}}\vspace{6pt}

Zeroth-order reinforcement learning for large language models suffers
from a fundamental noise amplification effect in the mismatched-data
setting: gradient variance scales as $1/\mu^2$, creating an irreducible
noise floor as the perturbation scale $\mu$ shrinks. The authors identify
the matched-data BOLT estimator as the antidote, derive ZO-BOLT, and
demonstrate 3--5$\times$ lower gradient variance and 40\% faster
convergence on GSM8K and HumanEval.

\begin{mdframed}[backgroundcolor=boxgray,linecolor=rulered,linewidth=1pt,
    innerleftmargin=6pt,innerrightmargin=6pt,
    innertopmargin=6pt,innerbottommargin=6pt]
\textbf{\small AT A GLANCE}\par\vspace{4pt}
\begin{description}[leftmargin=0pt,labelsep=4pt,itemsep=2pt,parsep=0pt,
    font=\small\bfseries]
  \item[Problem:] ZO RL for LLMs is unstable in practice; the reason
    is not understood theoretically.
  \item[Why it matters:] ZO methods are the only approach that enables
    policy-gradient RL without backpropagation, keeping memory at
    inference cost; instability blocks their practical use.
  \item[Method:] The authors prove a $1/\mu^2$ variance lower bound for
    mismatched-data ZO RL and derive ZO-BOLT, which shares rollout
    randomness between both perturbation evaluations to cancel the
    dominant noise term.
  \item[Result:] ZO-BOLT reduces gradient variance by 3--5$\times$
    and reduces forward passes to target accuracy by 40\% on
    GSM8K and HumanEval.
  \item[Evidence:] Gradient variance measurements; convergence curves on
    mathematical reasoning (GSM8K, MATH) and code generation
    (HumanEval).
\end{description}
\end{mdframed}

\section{The Big Picture}

Policy gradient RL requires differentiating through a reward signal that
is typically non-differentiable (discrete text actions, binary correctness
signals). Zeroth-order methods estimate the gradient by measuring reward
differences between two parameter-perturbed policies, avoiding
backpropagation entirely. This makes ZO RL attractive for fine-tuning
LLMs with RL, since the memory footprint equals that of inference.

Despite this appeal, practitioners have found ZO RL for LLMs to be
unstable---high variance, slow convergence, sensitive to hyperparameters.
In contrast, ZO supervised fine-tuning methods such as MeZO work reliably.
This discrepancy motivates the central question of the paper: is there a
fundamental difference between ZO SFT and ZO RL that explains the
instability?

\section{Why Existing Approaches Fall Short}

\textbf{Standard ZO RL} uses the two-point estimator:
\[
  \hat{g} = \frac{r(\theta + \mu z) - r(\theta - \mu z)}{2\mu} \cdot z
\]
where the two reward evaluations use independently sampled rollouts. The
mismatch between these rollouts is the fundamental problem.

\textbf{Smaller perturbation scale $\mu$} reduces the bias of the
estimator but, as the paper shows, the variance grows proportionally to
$1/\mu^2$---the two effects cannot be balanced simultaneously.

\textbf{More rollouts per update} reduces variance linearly but does not
address the $1/\mu^2$ structure, requiring a quadratic increase in rollouts
to halve the variance floor.

\section{Core Mechanism: The $1/\mu^2$ Lower Bound}

Let $r(\theta)$ be the stochastic reward with noise variance
$\sigma^2_\text{env}$. The paper shows that the variance of the standard
ZO RL estimator decomposes as:
\[
  \mathrm{Var}\!\left(\hat{g}\right)
  = \frac{\sigma^2_\text{env}}{2\mu^2}
  + O\!\left(\frac{\sigma^4_\text{env}}{\mu^4}\right)
  + \text{signal terms}
\]

The leading term $\sigma^2_\text{env}/(2\mu^2)$ grows without bound as
$\mu \to 0$, independent of sample size. This is the \emph{mismatch
penalty}: it arises specifically because the two evaluations $r(\theta +
\mu z)$ and $r(\theta - \mu z)$ use different rollout seeds, so their
noise does not cancel in the difference.

\section{The BOLT Solution}

Boltzmann Targeted SFT (BOLT) in the matched-data setting shares the
same rollout seed for both evaluations. Under matched randomness:
\[
  \tilde{g}_\text{BOLT} = \frac{r_s(\theta + \mu z) - r_s(\theta - \mu z)}{2\mu} \cdot z
\]
where both evaluations use rollout seed $s$. The noise now cancels:
\[
  r_s(\theta + \mu z) - r_s(\theta - \mu z)
  = \nabla_\theta r(\theta)^\top (2\mu z) + O(\mu^3)
\]
with no leading noise term. The variance reduces to:
\[
  \mathrm{Var}\!\left(\tilde{g}_\text{BOLT}\right) = O(1)
\]
independent of $\mu$, eliminating the $1/\mu^2$ amplification entirely.

\section{ZO-BOLT Algorithm}

ZO-BOLT instantiates matched-data perturbation for LLM RL fine-tuning:

\begin{enumerate}[itemsep=2pt,parsep=0pt]
  \item Sample a random direction $z$ in parameter space.
  \item Fix a rollout seed $s$.
  \item Evaluate the model under $\theta + \mu z$ with seed $s$ to
    obtain reward $r^+$.
  \item Evaluate the model under $\theta - \mu z$ with the same seed $s$
    to obtain reward $r^-$.
  \item Compute gradient estimate $\hat{g} = (r^+ - r^-)/(2\mu) \cdot z$.
  \item Update: $\theta \leftarrow \theta - \eta \hat{g}$.
\end{enumerate}

The matched-seed constraint adds zero computational overhead---it
requires only fixing the random seed before each rollout pair.

\section{What the Guarantee Says}

Theorem 1 (informal): For any ZO estimator in the mismatched-data
setting with environment noise $\sigma^2_\text{env} > 0$, the gradient
variance is at least $\sigma^2_\text{env}/(2\mu^2)$ for perturbation
scale $\mu$. No mismatched estimator can circumvent this bound.

Theorem 2 (informal): The ZO-BOLT matched-data estimator achieves
$O(1)$ variance, independent of $\mu$. Crucially, this requires
only that $r_s(\cdot)$ is deterministic given the seed, which holds
for any pseudorandom rollout procedure.

The bounds are non-asymptotic and hold for finite $d$ and $\mu > 0$.

\section{Experimental Findings}

\begin{itemize}[itemsep=2pt,parsep=0pt]
  \item \textbf{Gradient variance (GSM8K):} ZO-BOLT reduces gradient
    variance by $3.2\times$ over standard ZO RL at $\mu = 10^{-3}$,
    and $5.1\times$ at $\mu = 10^{-4}$.
  \item \textbf{Convergence (GSM8K):} ZO-BOLT reaches 70\% accuracy
    with 40\% fewer forward passes than standard ZO RL.
  \item \textbf{Convergence (HumanEval):} ZO-BOLT reaches 50\% Pass@1
    with 35\% fewer forward passes.
  \item \textbf{MATH:} ZO-BOLT final accuracy improves by 4.3\%
    over standard ZO RL with the same forward-pass budget.
\end{itemize}

\section{Ablations and Interpretation}

\textbf{Effect of $\mu$ on ZO-BOLT:} Unlike standard ZO RL, ZO-BOLT
performance is insensitive to $\mu$ across the range $10^{-4}$--$10^{-2}$,
confirming that the mismatch penalty has been eliminated.

\textbf{Seed granularity:} The paper evaluates fixing the seed per
rollout pair (default) vs.\ per batch. Per-pair seeding gives lower
variance, as expected from the theory.

\textbf{Comparison to first-order RL:} ZO-BOLT is $12\times$ slower
to converge than GRPO (a first-order RL method) on GSM8K, but uses
$8\times$ less memory---a favorable trade-off when GPU memory is the
bottleneck.

\section*{Reference}

Lianmin Chen, Junbin Qiu, Chenxing Wei, Yao Shu, Kun He.
\textbf{The Cost of Mismatch: Noise Amplification in Zeroth-Order
Reinforcement Learning.}
arXiv:2611.03260. NeurIPS 2026.
\url{https://arxiv.org/abs/2611.03260}

\end{document}
